\section{The scalar envelope and a barrier with nonmonotone speed}
\label{app:barrier-profiles}
\begin{proposition}[Exact envelope for a relaxed problem with constant $k=\log2$]
\label{prop:scalar-envelope}
Let $p\in C^2([0,\infty))$ satisfy $p(0)=0$, $p'>0$, $|p''|\leq p'$, and $p\leq p'$ on
$[0,\infty)$. Put $k=\log2$ and
\begin{equation}
 C(a)=\begin{cases}
 1+\displaystyle\frac1a\log\frac{k}{k(2-e^{-a})-a},
       &2a\leq k(2-e^{-a}),\\[5pt]
 1+\displaystyle\frac1a\log\frac{4a}{k(2-e^{-a})^2},
       &2a\geq k(2-e^{-a}).
 \end{cases}
 \label{eq:scalar-envelope-C}
\end{equation}
Then $C_{\SC}=\max_{a>0}C(a)$ has a unique maximizer, is strictly below two,
and
\[
 a\leq p^{-1}(ap'(a)/k)\leq C(a)a\leq C_{\SC}a.
\]
For the $C^2$ function $p$, put $\upsilon=p/p'$. The function $C(a)$ is
the exact envelope when $\upsilon$ is relaxed to any locally absolutely
continuous function $\upsilon:[0,\infty)\to[0,1]$ satisfying
\begin{equation}
 \upsilon(0)=0,\quad0<\upsilon(y)\leq1\quad(y>0),\qquad
 1-\upsilon\leq\upsilon'\leq1+\upsilon\quad\text{a.e.}
 \label{eq:upsilon-envelope}
\end{equation}
This exactness holds for the fixed constant $k$; it is not a claim about
the optimal arclength constant for products of copies of one smooth
barrier.
\end{proposition}
\begin{proof}
The derivative identity $\upsilon'=1-\upsilon p''/p'$ gives
\eqref{eq:upsilon-envelope}. Integration yields
\[
 1-e^{-a}\leq\upsilon(a)\leq\min(1,e^a-1),\qquad
 \frac{\upsilon(a)}a\leq\frac1k.
\]
For the last inequality, $(e^a-1)/a$ increases and $1/a$ decreases;
their crossing is at $a=k$. Thus the comparison point $a+t$,
defined by $p(a+t)=ap'(a)/k$, has $t\geq0$. Write
$u=\upsilon(a)$ and integrate $(\log p)'=1/\upsilon$:
\begin{equation}
 \int_0^t\frac{\dd q}{\upsilon(a+q)}=\log\frac{a}{ku}.
 \label{eq:envelope-time}
\end{equation}
The upper bound $\upsilon'\leq1+\upsilon$ and the cap $\upsilon\leq1$
give
$\upsilon(a+q)\leq m_u(q):=\min\{1,(1+u)e^q-1\}$.
The cap is reached at $q_u=\log(2/(1+u))$. For $t\leq q_u$,
\[
 \int_0^t\frac{\dd q}{m_u(q)}
 =\log\frac{(1+u)-e^{-t}}u.
\]
For $t\geq q_u$, the integral equals $\log((1+u)/(2u))+t-q_u$.
The right side of \eqref{eq:envelope-time} is fixed, and the actual
integral is at least this envelope integral. Solving the latter
therefore gives an upper bound on $t$:
\[
 t\leq\begin{cases}
 \log\dfrac{k}{k(1+u)-a},&2a\leq k(1+u),\\[3pt]
 \log\dfrac{4a}{k(1+u)^2},&2a\geq k(1+u).
 \end{cases}
\]
Both expressions decrease with $u$, and their values and first
derivatives agree at the switch. Substitute $u\geq1-e^{-a}$ to
obtain \eqref{eq:scalar-envelope-C}.

The two branches meet at exactly one positive point: $2a-k(2-e^{-a})$
has positive derivative $2-ke^{-a}$ and changes sign. It lies below $k$.
On the first branch, $C(a)<2$ is equivalent to
$2k(1-e^{-a})-a>0$. This strictly concave function vanishes at
$0$ and $k$, and the branch has $0<a<k$. On the second branch,
$C(a)<2$ is equivalent to
\[
 k(4e^a-4+e^{-a})-4a>0.
\]
Since $k>2/3$, the expression exceeds
$(2/3)(4e^a+e^{-a}-4-6a)$. The inequalities
$e^a\geq1+a+a^2/2+a^3/6$ and $e^{-a}\geq1-a$ bound the
parenthesis below by $1-3a+2a^2+(2/3)a^3$. Its minimum on
$[0,\infty)$ is $(16-5\sqrt{10})/3>0$, as differentiation shows.
Here $k>2/3$ follows, for example, from the first term of the
positive series for $\log2=2\operatorname{artanh}(1/3)$.

The two branches are continuously differentiable at their common
boundary. Their endpoint limits are $1/k$ as $a\downarrow0$ and
$1$ as $a\to\infty$. Moreover
$C(k)=1+k^{-1}\log(16/9)>1/k$, because
$\log(32/9)>1$ (use $e<3<32/9$).
Hence continuity proves a maximum at a finite positive argument,
and the strict pointwise bound implies $C_{\SC}<2$.
To prove uniqueness and to evaluate $C_{\SC}$ numerically, write $C=1+B/a$, where $B$ is
the logarithm on the appropriate branch. An interior stationary
point satisfies $aB'=B$, with
\[
 B'=\frac{1-ke^{-a}}{k(2-e^{-a})-a}\quad\hbox{or}\quad
 B'=\frac1a-\frac{2e^{-a}}{2-e^{-a}}.
\]
On the first branch, $B=-\log(2-e^{-a}-a/k)$ is strictly
convex, since its positive logarithm argument has second derivative
$-e^{-a}<0$. As $B(0)=0$, $aB'-B>0$ there. On the second,
\[
 B''=-\frac1{a^2}+\frac{4e^{-a}}{(2-e^{-a})^2}<0.
\]
Indeed the strict inequality is equivalent to
$2e^{a/2}-e^{-a/2}>2a$; the Taylor bounds
$e^t\geq1+t+t^2/2$ and $e^{-t}\leq1-t+t^2/2$ for $t\geq0$
give a difference at least $1-a/2+a^2/8>0$.
Thus $aB'-B$ strictly decreases on the second branch, starts
positive by continuous differentiability at the switch, and tends
to $-\infty$. It has one zero, which is the unique global maximizer.
Numerical evaluation gives $a\approx0.65011438345$ and
$C_{\SC}\approx1.831856423$; these decimals are not used in the
bound or in its attainment proof.

Finally, for fixed $a$ take $\upsilon(y)=1-e^{-y}$ until $a$,
then follow $\upsilon'=1+\upsilon$ until reaching one, and
remain one thereafter. This continuous, piecewise differentiable
function obeys \eqref{eq:upsilon-envelope} almost everywhere and
attains every envelope equality used at $a$. Reconstruct $p$ by
$(\log p)'=1/\upsilon$, with normalization $p'(0)=1$;
the behavior $\upsilon(y)=y+O(y^2)$ gives a well-defined positive
solution with $p(0)=0$. This proves exactness for the stated relaxed
problem. Likewise, $k$ is the largest constant valid for all
admissible functions $\upsilon$, since the admissible function
$\upsilon=\min(1,e^y-1)$ attains $\inf_y y/\upsilon(y)=k$. These two
extremal functions need not come from one fixed smooth scalar barrier.
\end{proof}

\subsection{A smooth barrier whose central speed is not monotone}
\label{app:nonmonotone-construction}
For $A>4$ define, for all real $y$,
\[
 P_A(y)=\frac{1+2e^{-A}\cosh y}{1+2e^{-A}},\quad
 X_A(y)=\int_0^y\frac{\dd u}{P_A(u)},\quad
 Z_A(y)=\int_0^y P_A(u)\dd u.
\]
The map $X_A$ is a smooth increasing bijection from $\R$ onto
a finite interval $(-b_A,b_A)$: $P_A(y)\geq
e^{y-A}/(1+2e^{-A})$ for $y\geq0$, and $P_A$ is even.
Define $f_A$ by $f_A'(X_A(y))=Z_A(y)$ and $f_A(0)=0$.
All these functions are smooth, so no piecewise smoothing or boundary
extension is needed.
Differentiation gives
\begin{equation}
 f_A''(X_A(y))=P_A(y)^2,
 \quad\frac{|f_A'''|}{(f_A'')^{3/2}}=
        2\frac{|P_A'|}{P_A}<2,
 \quad V_A(y):=\frac{f_A'}{\sqrt{f_A''}}
 =\frac{y+2e^{-A}\sinh y}{1+2e^{-A}\cosh y}.
 \label{eq:nonmonotone-barrier}
\end{equation}
In particular, $y$ is the exact metric coordinate of $f_A$, and $0$ is
its analytic center. Since $V_A(y)\to\pm1$ as $y\to\pm\infty$,
the identity $\frac{\dd}{\dd y}f_A(X_A(y))=V_A(y)$ proves blow-up
at both endpoints. Also $Z_A(y)\to\pm\infty$,
so $f_A'$ takes every real value.

The exact parameter of $f_A$ is $\nu_A=\sup_y V_A(y)^2$.
For $0\leq y\leq A$, $V_A(y)\leq y+1\leq A+1$.
For $y=A+u\geq A$, use $2e^{-A}\sinh y\leq2e^{-A}\cosh y$
and $2e^{-A}\cosh y\geq e^u$ to get
\[
 V_A(y)\leq1+\frac{A+u-1}{1+e^u}\leq A+1;
\]
the final inequality follows from $u-1\leq Ae^u$ for $A>4$.
Oddness handles negative $y$. At $y=A/2$, the denominator in
\eqref{eq:nonmonotone-barrier} is
$1+e^{-A/2}+e^{-3A/2}<3/2$, so $V_A(A/2)>A/3$.
Thus $A^2/9<\nu_A\leq(A+1)^2$. The speed $V_A$ is not monotone:
it exceeds one at $A/2$ but tends to one at the endpoint.

Fix $r$ and $0<\delta<1/2$. Uniformly for
$y\in[\delta A,(1-\delta)A]$ as $A\to\infty$,
$P_A(y)=1+o(1)$ and $Z_A(y)=y+o(1)$.
Consequently, while the metric coordinate crosses this interval, the
logarithmic central-path parameter increases by
\[
 \log Z_A((1-\delta)A)-\log Z_A(\delta A)
 =\log\frac{1-\delta}{\delta}+o(1).
\]
Choose $B>\log((1-\delta)/\delta)$ and weights with
$\log w_i=B(r-i)$, and stop the central path at
$s_1=\log Z_A((1-\delta)A)$. For large $A$, the $r$ parameter windows,
shifted by the weights, are disjoint, and in each one metric coordinate
advances by $(1-2\delta)A$. Therefore the central arc has length at
least $r(1-2\delta)A$.

At the final parameter, $Z_A(y_i)=Z_A((1-\delta)A)e^{B(r-i)}$.
For $y\geq\delta A$ and large $A$,
$Z_A(y)\geq e^{y-A}/3$ and
$Z_A((1-\delta)A)<A$. Thus
$y_i\leq A+\log(3A)+B(r-i)$, so the endpoint distance is at most
$\sqrt r[A+\log(3A)+Br]$.
Let $A\to\infty$ with $r,\delta,B$ fixed, then let
$\delta\downarrow0$. The ratio approaches $\sqrt r$, proving
\cref{prop:nonmonotone}. This does not contradict comparisons that
depend on the barrier parameter, such as \eqref{eq:NN-comparison}: the
scalar parameter, and hence the product parameter $r\nu_A$, diverges in
this construction.
